\documentclass[aspectratio=169]{beamer}
\usetheme{veit}
\title{Abstractions! How do I even?}
\date{\today}
\author{Veit Heller}
\institute{EnthusiastiCon 2019}
\setbeamertemplate{section page}
{
    \begingroup
    \usebeamerfont{section title}\insertsection\par
    \endgroup
}
\AtBeginSection[]{%
  \begin{frame}[plain]
    \usebeamertemplate{section page}
  \end{frame}
}
\begin{document}
  \maketitle
  \begin{frame}{Agenda}
    \begin{itemize}
      \item Motivation!
      \item Michel Foucault—The Order of Things
      \item Robert M. Pirsig—Zen and the Art of Motorcycle Maintenance
      \item Summary!
    \end{itemize}
  \end{frame}
  \section{Motivation!}
  \begin{frame}{Technophilosophy?}
    We’re going to talk a lot about philosophy. None of it will be about the
    branches of philosophy that seem most notably intertwined with our
    profession: technophilosophy.
  \end{frame}
  \begin{frame}{Abstractions}
    We build abstractions, always.
  \end{frame}
  \begin{frame}{Abstractions}
    So do philosophers; they have a headstart of just a few millenia.
  \end{frame}
  \begin{frame}{But I don’t have a PhD!}
    Me neither, and English isn’t my first language.

    The beauty about a lot of philosophy is that it is very simple, because it
    deals with very simple things—but not easy things!
  \end{frame}
  \section{Michel Foucault—The Order of Things}
  \begin{frame}{Disclaimer}
    This book is not about abstractions, but about “episteme”, which are roughly
    the fundamental ideological axioms of all of knowledge—think paradigms, but
    without knowing that you adhere to them. This doesn’t make any sense
    without context, and the talk isn’t really about it anyway.

    What it is about is how to impose order through abstraction.
  \end{frame}
  \begin{frame}{Jorge Luis Borges}
    “[\ldots] animals are divided into: (a) belonging to the Emperor, (b)
    embalmed, (c) tame, (d) sucking pigs, (e) sirens, (f) fabulous, (g) stray
    dogs, (h) included in the present classification, (i) frenzied, (j)
    innumerable, (k) drawn with a very fine camelhair brush, (l) \textit{et
    cetera}, (m) having just broken the water pitcher, (n) that from a long
    way off look like flies [\ldots]”

    — Michel Foucault, citing Jorge Luis Borges, in “The Order of Things”
  \end{frame}
  \begin{frame}{Jose Luis Borges, Programmer}
    \includegraphics[width=12cm]{traceback}
  \end{frame}
  \begin{frame}{Classifications}
    Classifications are defined by the space between them, not by their label.
  \end{frame}
  \begin{frame}{Why classifications?}
    If the classifications are intuitive, the abstraction will be intuitive as
    well.
  \end{frame}
  \begin{frame}{Why classification?}
    Side note: it’s impossible to abstract without classifying, because
    abstractions always add things or leave them out.
  \end{frame}
  \begin{frame}{Why no classification?}
    Classification and order is hierarchy—don’t let yourself be sucked into
    using hierarchies as a lens to see the world.
  \end{frame}
  \section{Robert M. Pirsig—Zen and the Art of Motorcycle Maintenance}
  \begin{frame}{Disclaimer}
    This book is good, and reading it is very addicting!

    It’s about a father-son motorcycle trip, a life changed by mental illness,
    and the nature of “Quality”.
  \end{frame}
  \begin{frame}{The Three Things}
    There are three things in all dialogues: the subject, the object, and
    “Quality”, which happens between the two.
  \end{frame}
  \begin{frame}{The Three Things}
    Quality is always between the observer and the observed.
  \end{frame}
  \begin{frame}{The Three Things}
    Abstractions are always subjective. There is always a use case.
  \end{frame}
  \begin{frame}{Objectivity in a subjective world}
    If everything is subjective, how do we empower opinionated users?
  \end{frame}
  \begin{frame}{The curious case of Git}
    Git is the ultimate leaky abstraction.
  \end{frame}
  \begin{frame}{The curious case of Git}
    It gives you control over the “porcelain”, or user-facing facilities, but
    also over the “plumbing”, and it does so gracefully.

    \small{%
      I can \texttt{git commit}, but I can just as well \texttt{git commit-tree
      <tree> -p <parent-commit>}.
    }
  \end{frame}
  \begin{frame}{Do abstractions have to hide their details?}
    Only if you consider your abstractions to be good enough for everyone.
  \end{frame}
  \section{Summary!}
  \begin{frame}{Why do I have to think about these things?}
    Thinking about how to write abstractions will make you reflect about your
    personal æsthetics and whether what you build actually holds up to your
    standards.
  \end{frame}
  \begin{frame}{Why do I have to think about these things?}
    Reading philosophy will inspire you, or at least build a vocabulary that
    heps you articulate your æsthetics.
  \end{frame}
  \begin{frame}{Techniques for writing better abstractions, summarized}
    A personal guidebook.
    \begin{itemize}
      \item Don’t fear taking a hard stance if it leads to better
            classifications and the right abstractions. But be gentle to
            other people if they disagree.
      \item There should be no overlap between different function sets in
            your API. Remember: it’s about clear boundaries.
      \item If the order is intuitive, the abstraction will be intuitive as
            well. But please write documentation anyway.
      \item If you’re unsure about anything, talk to someone else. Bee tee
            dubbs: I’m always happy to chat about these things.
    \end{itemize}
  \end{frame}
  \begin{frame}{A Commentated Reading List I}
    \begin{itemize}
      \item Michel Foucault—The Order of Things: not about abstractions per se,
            but nonetheless exciting!
      \item Robert M. Pirsig—Zen and the Art of Motorcycle Maintenance: one of
            the best books I’ve ever read, stock-full of marvelous writing,
            storytelling, and compelling philosophical ideas.
      \item Christopher Alexander—Notes on the Synthesis of Form: provides a
            beautiful system of judging tradeoffs and implementations of ideas.
      \item Zachary Tellman—Elements of Clojure: talks about names and idioms
            rather than how to build anything in particular, and is thus more
            useful than any other programming book I’ve read; not just for
            Clojure programmers.
      \item Guy Steele—Growing a Language: both a paper and a talk. It explores
            building technical abstractions from first principles, and is
            removed enough from technology to be generally useful.
    \end{itemize}
  \end{frame}
  \begin{frame}
    \begin{itemize}
      \item Douglas R. Hofstaedter—Gödel, Escher, Bach: An Eternal Golden Braid:
            a weird and beautiful voyage into meaning, self-reference, and
            recursion.
      \item The Git Book, Chapter 10.1—Git Internals - Plumbing and Porcelain:
            is the start of a journey into weird Git internals that somehow ends
            up being extremely empowering and beautiful.
    \end{itemize}
  \end{frame}
  \begin{frame}{The End}
    \Huge Thank you!
    \linebreak
    \linebreak
    \tiny Slides at \texttt{https://github.com/hellerve/talks}
  \end{frame}
\end{document}
